\section{Introducci\'on}

\subsection{Contexto y motivaci\'on}

La caracterizaci\'on geot\'ecnica del suelo es una etapa fundamental en el dise\~no de
cimentaciones, obras civiles, estructuras sometidas a cargas din\'amicas y, en general,
en cualquier proyecto que requiera conocer el comportamiento mec\'anico del
terreno~\cite{Kramer1996,Mitchell2005}. El par\'ametro de referencia para el an\'alisis de
respuesta s\'ismica de sitio es el perfil de velocidad de onda de corte $V_S(z)$, que se
vincula directamente con el m\'odulo de corte a peque\~nas deformaciones $G_{\max}$
mediante $V_S=\sqrt{G_{\max}/\rho}$.

En la pr\'actica profesional local, la caracterizaci\'on se apoya casi exclusivamente en
ensayos invasivos como el \emph{Standard Penetration Test} (SPT) y el \emph{Cone
Penetration Test} (CPT)~\cite{ASTM}. Estos ensayos entregan informaci\'on valiosa y
ampliamente aceptada, pero comparten un conjunto de limitaciones: requieren
intervenci\'on directa del terreno, implican tiempos de ejecuci\'on considerables, tienen
costos operativos elevados, presentan dificultades log\'isticas en zonas de acceso
complejo, y ---por su naturaleza puntual--- entregan informaci\'on de un volumen de
suelo muy reducido en torno al punto de sondeo.

Estas limitaciones motivan el estudio de m\'etodos complementarios de menor impacto.
Entre las t\'ecnicas geof\'isicas aplicables se destacan la Tomograf\'ia de Refracci\'on
S\'ismica (SRT) y el An\'alisis Multicanal de Ondas Superficiales
(MASW)~\cite{Park1999,Xia1999,Foti2014}. El MASW explota la \emph{dispersi\'on
geom\'etrica} de las ondas Rayleigh: componentes de distinta longitud de onda muestrean
distintas profundidades del semiespacio, de modo que la dependencia de la velocidad de
fase con la frecuencia codifica la estratificaci\'on del subsuelo. Es un m\'etodo no
invasivo, r\'apido, de costo relativamente bajo, y que integra un volumen de suelo
lateralmente mucho m\'as representativo que un sondeo puntual.

\subsection{El problema instrumental}

La adopci\'on del MASW en el medio local est\'a limitada menos por la teor\'ia ---que est\'a
madura y bien documentada--- que por el acceso al instrumental. Un sismógrafo
multicanal comercial con sus ge\'ofonos representa una inversi\'on que dif\'icilmente se
justifica fuera de una empresa especializada, y el cableado multicanal tradicional
impone una log\'istica de despliegue pesada en campo.

Desde el punto de vista electr\'onico, construir un sistema propio no es trivial. El
ge\'ofono de bobina m\'ovil entrega se\~nales del orden de las decenas o centenas de
microvoltios, con un ancho de banda \'util que arranca por debajo de su frecuencia
natural. Esto exige ganancias altas, y en cadenas de acondicionamiento multietapa las
tensiones de \emph{offset} de los amplificadores y de las referencias se acumulan
etapa a etapa, consumiendo rango din\'amico y, con ganancia elevada, llegando a saturar
etapas intermedias antes de que la se\~nal s\'ismica alcance el conversor. A esto se suma
el requisito de sincronizaci\'on: en un arreglo multicanal, los desfasajes temporales
entre canales se traducen directamente en errores de velocidad de fase, es decir, en
errores del perfil $V_S$ estimado.

\subsection{Objetivo del proyecto}

El objetivo general es \textbf{dise\~nar, implementar y evaluar experimentalmente un
sistema electr\'onico multicanal para la generaci\'on y captaci\'on de ondas s\'ismicas},
orientado a la caracterizaci\'on geot\'ecnica del suelo en profundidades de hasta
aproximadamente \SI{50}{\meter}. El sistema comprende la selecci\'on e implementaci\'on
del generador de ondas s\'ismicas, el dise\~no de nodos ge\'ofono aut\'onomos del tipo IoT,
su organizaci\'on en configuraciones matriciales de medici\'on, y el procesamiento digital
de las se\~nales adquiridas para estimar perfiles de velocidad y par\'ametros el\'asticos
de las subcapas del suelo.

Los objetivos espec\'ificos, derivados de las etapas de ejecuci\'on aprobadas en la
propuesta, son:

\begin{enumerate}[itemsep=0.15em]
    \item Estudiar los principios de propagaci\'on de ondas s\'ismicas en suelos y revisar
    el estado del arte sobre SRT, MASW, ge\'ofonos, adquisici\'on multicanal y sistemas
    IoT aplicados a instrumentaci\'on de campo.
    \item Seleccionar el generador de ondas s\'ismicas, los ge\'ofonos y los componentes
    electr\'onicos principales para la captaci\'on, acondicionamiento, digitalizaci\'on y
    transmisi\'on de las se\~nales.
    \item Dise\~nar los circuitos de acondicionamiento, alimentaci\'on, sincronizaci\'on y
    registro de datos de los nodos ge\'ofono aut\'onomos.
    \item Implementar el prototipo electr\'onico multicanal y definir las topolog\'ias de
    arreglo a emplear en las pruebas de campo.
    \item Desarrollar los algoritmos de procesamiento digital, incluyendo filtrado,
    detecci\'on de llegadas, an\'alisis de dispersi\'on e inversi\'on para estimaci\'on de
    perfiles de velocidad.
    \item Integrar el sistema, realizar ensayos experimentales, comparar los resultados
    con el ensayo SPT, analizar el desempe\~no y redactar la memoria final.
\end{enumerate}

\subsection{Estado de avance y alcance de este documento}

Al momento de esta Primera Presentaci\'on se encuentran completadas las etapas 1 a 5 en
su mayor parte, con la etapa 6 iniciada: existe un prototipo funcional de varios nodos
que captura, transmite y almacena datos s\'ismicos reales, se realizaron campa\~nas de
medici\'on en campo, y el \emph{pipeline} de procesamiento MASW produce curvas de
dispersi\'on e inversiones de perfil $V_S$. Lo que resta es principalmente la
consolidaci\'on del hardware en una PCB propia, la construcci\'on de una fuente s\'ismica
repetible, y la validaci\'on cuantitativa contra un ensayo de referencia.

Un resultado parcial del trabajo ---el \emph{front-end} anal\'ogico asistido digitalmente
con auto-calibraci\'on de \emph{offset}--- fue redactado como comunicaci\'on cient\'ifica y
sometido al congreso URUCON~2026, lo que aport\'o al proyecto un ciclo de revisi\'on
externa sobre una parte central de la instrumentaci\'on.

El resto del documento se organiza como sigue. La Secci\'on~\ref{sec:fundamentos}
desarrolla los fundamentos te\'oricos de propagaci\'on de ondas superficiales. La
Secci\'on~\ref{sec:estado-arte} presenta el estado del arte en m\'etodos, instrumentaci\'on
y herramientas de procesamiento. La Secci\'on~\ref{sec:arquitectura} describe la
arquitectura del sistema desarrollado. La Secci\'on~\ref{sec:metodologia} detalla la
metodolog\'ia experimental y los procedimientos de validaci\'on. La
Secci\'on~\ref{sec:resultados} presenta los resultados obtenidos. La
Secci\'on~\ref{sec:futuros} enumera los trabajos futuros y la
Secci\'on~\ref{sec:conclusiones} cierra con las conclusiones parciales.

\nota{Para la versi\'on de 15 p\'aginas, esta introducci\'on deber\'ia comprimirse a
aproximadamente una p\'agina: conservar \S1.1 y \S1.3 (motivaci\'on y objetivos),
resumir \S1.2 en un p\'arrafo y eliminar el mapa de secciones.}
